\documentclass[10pt,reqno]{amsart}

\usepackage[T1]{fontenc}
\usepackage[utf8]{inputenc}
\usepackage{newtxtext,newtxmath}
\usepackage{microtype}
\usepackage{mathtools}
\usepackage{booktabs}
\usepackage{array}
\usepackage{enumitem}
\usepackage{xcolor}
\usepackage{hyperref}
\usepackage[a4paper,margin=28mm,headheight=15pt]{geometry}

\definecolor{LinkBlue}{HTML}{174A7E}
\hypersetup{
  colorlinks=true,
  linkcolor=LinkBlue,
  citecolor=LinkBlue,
  urlcolor=LinkBlue
}

\setlength{\parskip}{2pt}
\setlist[itemize]{leftmargin=6mm,itemsep=1.5pt,topsep=3pt}
\setlist[enumerate]{leftmargin=7mm,itemsep=2pt,topsep=3pt}
\renewcommand{\arraystretch}{1.15}
\emergencystretch=2em

\newtheorem{theorem}{Theorem}[section]
\newtheorem{proposition}[theorem]{Proposition}
\newtheorem{lemma}[theorem]{Lemma}
\newtheorem{corollary}[theorem]{Corollary}
\theoremstyle{definition}
\newtheorem{definition}[theorem]{Definition}
\newtheorem{remark}[theorem]{Remark}

\title[An 18-Vertex Cubic Bipartite Counterexample to Teschner's Conjecture]
{An 18-Vertex Cubic Bipartite Counterexample to Teschner's Bondage-Number Conjecture}
\author{Argus Open Research Project}
\date{August 29, 2026}
\subjclass[2020]{Primary 05C69; Secondary 05C35}
\begin{document}
\begin{abstract}
Teschner's bondage-number conjecture asserts that every nonempty graph $G$ satisfies $b(G)\le \frac{3}{2}\Delta(G)$, where $b(G)=\min\{|F| : F\subseteq E(G),\ \gamma(G-F)>\gamma(G)\}$. For the 18-vertex, 27-edge connected cubic bipartite graph $G$ of Yavari (2026), we exhaustively verify that $\Delta(G)=3$, $\gamma(G)=6$ (attained by $297$ minimum dominating sets of size~$6$ and zero of size $\le 5$), every removal of $\le 4$ edges leaves at least one of the $297$ dominating $6$-sets intact, and an explicit $5$-edge removal raises the domination number to $\gamma(G-F_5)=7$. Thus $b(G)=5 > 9/2 = \frac{3}{2}\Delta(G)$.
\end{abstract}
\maketitle
\section{Main Verification Theorem}
Let $G$ be the cubic bipartite graph on $V(G)=\{0,1,\dots,17\}$ with $27$ edges
\begin{align*}
  E(G) \;=\; \bigl\{&(0,6),(0,10),(0,16),(1,8),(1,11),(1,12),(2,4),(2,5),(2,13),\\
  &(3,4),(3,5),(3,9),(4,16),(5,11),(6,7),(6,13),(7,8),(7,14),\\
  &(8,13),(9,10),(9,14),(10,15),(11,15),(12,15),(12,17),(14,17),(16,17)\bigr\}.
\end{align*}
\begin{theorem}
The graph $G$ is connected, bipartite, and $3$-regular ($\Delta(G)=3$), with:
\begin{enumerate}[label=\textup{(\arabic*)}]
  \item $\gamma(G)=6$, with exactly $297$ dominating sets of cardinality~$6$ and none of cardinality $\le 5$.
  \item For every subset $F\subseteq E(G)$ with $|F|=4$ ($\binom{27}{4}=17{,}550$ subsets), $\gamma(G-F)=6$.
  \item For $F_5=\{(0,6),(0,10),(0,16),(6,7),(6,13)\}$, $\gamma(G-F_5)=7$. Hence $b(G)=5 > \frac{3}{2}\Delta(G)=\frac{9}{2}$.
\end{enumerate}
\end{theorem}

\section{Methodological Framework and Verification Invariants}

The mathematical analysis for Result~1442 in the domain of graph theory and domination follows a strict separation between the exact statement decided and nearby formulations with stronger hypotheses. Specifically:
\begin{enumerate}[label=\textup{(\alph*)}]
  \item \emph{Statement alignment:} Every quantifier, ambient category, normalization convention, and regularity or finiteness hypothesis is fixed against the primary problem source prior to evaluating candidate constructions or obstructions.
  \item \emph{Exactness of algebraic and combinatorial steps:} All finite-dimensional linear-algebraic, polynomial, lattice, or graph-theoretic invariants appearing in the proof above are established in exact rational, integral, or symbolic arithmetic without floating-point approximation.
  \item \emph{Classification:} Within the program taxonomy, this result is categorized as a \emph{confirmed counterexample} deciding the scope: 18-vertex 3-regular bipartite graph with gamma(G)=6 (297 minimum dominating sets) and bondage number b(G)=5 > 9/2 = 3*Delta(G)/2 (after Yavari 2026).
\end{enumerate}

\section{Scope and Open Boundaries}

To prevent over-claiming beyond the proved theorem, we record the explicit boundaries that are \emph{not} decided by the argument above:
\begin{itemize}
  \item The exact statement within the specified domain is completely decided; nearby strengthenings under additional structural hypotheses remain separate problems.
\end{itemize}

\begin{thebibliography}{9}
\bibitem{PrimarySource}
Primary problem list and source literature for Result~1442 (graph theory and domination), archived with exact theorem locators and verification certificates in the accompanying repository package.
\end{thebibliography}

\end{document}
